\subsection{\fid{F64}: weight grading does not preserve the actual flag convolution}
\label{sec:weightextraction}

The nonconstant torus coefficient of \fid{F62} suggests passing
to weights or perverse cohomology before seeking positive
coefficients. The following tests these specified operations on
the same actual standard flag kernels. It also records the
proper positive construction obtained by choosing a different
generator. No assertion about every filtered or nonproper
enhancement is made.

Retain \(k=\mathbb F_q\), the coefficient field \(E\),
\(X=\mathbb P^1_k\), \(Y=X\times X\), and the inclusions
\(i:\Delta\hookrightarrow Y\), \(j:O=Y\setminus\Delta
\hookrightarrow Y\) from \fid{F62}. Put
\[
 \mathcal I=i_*E_X,\qquad
 \mathcal B=E_Y,\qquad
 \mathcal H=j_!E_O.
\]
Convolution is the actual pull--tensor--compact-push operation
\eqref{eq:flagkernelconvolution}, with unit \(\mathcal I\).

For a mixed constructible complex \(\mathcal A\), define
the additive splitting operation
\begin{equation}\label{eq:weightgradedoperation}
 \mathsf G(\mathcal A)=
 \bigoplus_{r,w}
       \operatorname{gr}^W_w\,{}^pH^r(\mathcal A)[-r].
\end{equation}
The finite perverse weight filtration is canonical and strictly
functorial \cite[Theorem~5.3.5 and its proof through
Lemma~5.3.6, pp.~136--137]{BBD}. We retain the original Weil
structures on the graded pieces. On morphisms,
\eqref{eq:weightgradedoperation} applies perverse cohomology,
the functorial weight graded map, and finite direct sum.
Thus it is a defined additive operation; no triangulated
exactness or compatibility with convolution is assumed.

It preserves the Grothendieck class and every finite-extension
stalk Euler trace. Indeed, perverse cohomology with its shifts
reconstructs the derived Euler class, and each finite weight
filtration reconstructs its perverse class. This is only an
additive Euler assertion.

Open restriction is perverse \(t\)-exact: it is the smooth
relative-dimension-zero case of the discussion immediately
before \cite[Proposition~4.2.5, pp.~108--109]{BBD}.
It preserves both weight bounds, since \(j^*=j^!\) and the two
bounds in \cite[5.1.14, p.~128]{BBD} apply. Restricting a
weight filtration therefore gives pure graded pieces of the
same weights or zero. Uniqueness identifies this with the
weight filtration of the restriction. Consequently
\begin{equation}\label{eq:weightgradedopenrestriction}
             j^*\mathsf G(\mathcal A)
                  \simeq\mathsf G(j^*\mathcal A).
\end{equation}
This is not a claim about purity under arbitrary closed
restriction.

\begin{finding}[\fid{F64}, a complete extraction test]
\label{f:weightextraction}
For these actual standard kernels,
\begin{equation}\label{eq:weightgradedH}
                  \mathsf G(\mathcal H)
                     =\mathcal B\oplus\mathcal I[-1],
\end{equation}
but
\begin{equation}\label{eq:weightgradednotmonoidal}
 \mathsf G(\mathcal H\star\mathcal H)
       \not\simeq
 \mathsf G(\mathcal H)\star\mathsf G(\mathcal H).
\end{equation}
Splitting into perverse cohomology alone, or into ordinary
cohomology alone, likewise fails to commute with this product.

Moreover there is no exact strong-convolution functor from
a triangulated kernel subcategory containing
\(\mathcal H,\mathcal I\), the required scalar shifts and
localization triangles, to the same standard kernel category,
which sends \(\mathcal I\) to \(\mathcal I\), commutes with
Tate twists and cohomological shifts, and sends
\(\mathcal H\) to \(\mathcal B\).
\end{finding}

\begin{proof}
The localization triangle is
\begin{equation}\label{eq:weightflaglocalization}
 \mathcal H\longrightarrow\mathcal B
       \longrightarrow\mathcal I\longrightarrow\mathcal H[1].
\end{equation}
Put \(P=\mathcal B[2]\) and \(Q=\mathcal I[1]\).
They are perverse and pure of weights \(2\) and \(1\):
they are the shifted constant sheaves on smooth varieties,
with closed pushforward for \(Q\)
\cite[\S4.0, p.~102; 5.1.8, pp.~126--127]{BBD}.
Shift \eqref{eq:weightflaglocalization} by \(2\) and take
perverse cohomology. Its only nonzero part is
\begin{equation}\label{eq:openkernelweightfiltration}
 0\longrightarrow Q\longrightarrow\mathcal H[2]
     \longrightarrow P\longrightarrow0.
\end{equation}
In particular \(\mathcal H[2]\) is perverse. This conclusion
uses the triangle directly, without an extra affine
\(t\)-exactness premise. Since the end terms have weights
\(1\) and \(2\), uniqueness identifies
\eqref{eq:openkernelweightfiltration} with the weight
filtration:
\[
 W_1(\mathcal H[2])=Q,\qquad
 \operatorname{gr}^W_2(\mathcal H[2])=P.
\]
Equation \eqref{eq:weightgradedH} follows. Also
\(\mathsf G(\mathcal B)=\mathcal B\) and
\(\mathsf G(\mathcal I)=\mathcal I\).

The unit relation and the proper middle-point calculation give
\begin{equation}\label{eq:closedflagkernelrelation}
 \begin{split}
 \mathcal B\star\mathcal I
   &=\mathcal I\star\mathcal B=\mathcal B,\\
 \mathcal B\star\mathcal B
   &\simeq\mathcal B\otimes R\Gamma(\mathbb P^1,E)
     \simeq\mathcal B\oplus\mathcal B(-1)[-2].
 \end{split}
\end{equation}
These are actual geometric identifications: the middle point
is an independent proper \(\mathbb P^1\), whose cohomology is
\(E\) in degree zero and \(E(-1)\) in degree two.
Distribution over finite direct sums now yields
\begin{equation}\label{eq:productofweightgradedH}
 \mathsf G(\mathcal H)\star\mathsf G(\mathcal H)
  \simeq
   \mathcal B\oplus\mathcal B[-1]^{\oplus2}
      \oplus\mathcal B(-1)[-2]\oplus\mathcal I[-2].
\end{equation}
Its restriction to \(O\) has nonzero ordinary degree-zero
cohomology \(E_O\).

On the other hand, \fid{F62} identifies
\(j^*(\mathcal H\star\mathcal H)=\mathcal K\), with
\[
 H^1(\mathcal K)=E_O,\qquad
 H^2(\mathcal K)=E_O(-1),
\]
and no other ordinary cohomology sheaves. They are lisse on
the smooth two-dimensional \(O\). The perverse \(t\)-structure
on complexes of lisse sheaves is the ordinary one shifted by
two; hence
\begin{equation}\label{eq:torusperverseindices}
 {}^pH^3(\mathcal K)=E_O[2],\qquad
 {}^pH^4(\mathcal K)=E_O(-1)[2],
\end{equation}
and all other perverse cohomology vanishes. Equivalently these
indices follow by applying perverse cohomology to the Chern
triangle \eqref{eq:torusCherntriangle}; its two end terms
are in perverse degrees four and two. The perverse sheaves
in \eqref{eq:torusperverseindices} are already pure, of weights
\(2\) and \(4\). By
\eqref{eq:weightgradedopenrestriction},
\begin{equation}\label{eq:weightgradedproductrestriction}
 j^*\mathsf G(\mathcal H\star\mathcal H)
       \simeq E_O[-1]\oplus E_O(-1)[-2].
\end{equation}
Its ordinary degree-zero cohomology is zero. This proves
\eqref{eq:weightgradednotmonoidal}.

For perverse-cohomology splitting put
\[
       \mathsf P(\mathcal A)=
            \bigoplus_r{}^pH^r(\mathcal A)[-r].
\]
Since \(\mathcal H[2]\) is perverse,
\(\mathsf P(\mathcal H)=\mathcal H\). But
\eqref{eq:torusperverseindices} gives
\[
 \begin{split}
 j^*\mathsf P(\mathcal H\star\mathcal H)
   &=E_O[-1]\oplus E_O(-1)[-2],\\
 j^*(\mathsf P(\mathcal H)\star\mathsf P(\mathcal H))
   &=\mathcal K.
 \end{split}
\]
The nonzero Chern class in \fid{F62} proves that these are
not isomorphic. The same proof applies to the ordinary
splitting \(\mathsf S(\mathcal A)=\bigoplus_r
H^r(\mathcal A)[-r]\), since \(\mathcal H\) is an ordinary
sheaf. In addition, \fid{F62} gives total cohomology dimension
two for \(R\Gamma(O,\mathcal K)\), but four for
\(R\Gamma(O,\mathsf S(\mathcal K))\). Thus this splitting
also fails to commute with that actual relative compact
pushforward followed by ordinary base cohomology.

It remains to prove the exact-functor assertion.
Let \(L=[E(-1)[-2]]\) denote the scalar coefficient class.
The two actual localization triangles of \fid{F62} give
\begin{equation}\label{eq:openflagKzerorelation}
       [\mathcal H]^2=(L-1)[\mathcal H]+L[\mathcal I].
\end{equation}
An exact functor induces an additive map on \(K_0\), and
strong monoidality preserves products. The stated scalar,
unit and generator comparisons would force
\[
        [\mathcal B]^2=(L-1)[\mathcal B]+L[\mathcal I].
\]
The actual proper relation
\eqref{eq:closedflagkernelrelation} instead gives
\([\mathcal B]^2=(L+1)[\mathcal B]\). Subtracting yields
\(2[\mathcal B]=L[\mathcal I]\). Restriction to \(O\)
then gives \(2[E_O]=0\), impossible because ordinary
generic Euler rank is an integer-valued homomorphism and
takes this class to \(2\).
\end{proof}

\paragraph{What the counterexample preserves and loses.}
On the open stratum, the difference between
\eqref{eq:productofweightgradedH} and
\eqref{eq:weightgradedproductrestriction} is
\(E_O\oplus E_O[-1]\). It is a nonzero derived object with
zero Euler class. Thus failure of the operation comparison
is compatible with equality of every trace function. The
proof uses the canonical weight filtration, not a chosen
splitting of an arithmetic Frobenius module.

The quotient \(\mathcal H[2]\to\mathcal B[2]\) in
\eqref{eq:openkernelweightfiltration} is the familiar
intermediate-extension quotient:
\[
                        j_{!*}E_O[2]=E_Y[2].
\]
The constant irreducible local system extends to the smooth
irreducible \(Y\), so this identity follows from
\cite[Lemma~4.3.2 and proof, pp.~112--113]{BBD}.
It is a canonical proper pure candidate. The exact-functor
part of the finding explains why selecting it changes the
generator and its relation. Without exactness, preservation
of the additive Euler relation is not automatic; without
the specified scalar or unit comparison, that proof does
not apply. It does not exclude a changed algebra presentation.

\paragraph{The positive proper construction that does exist.}
For the closed kernel \(\mathcal B\), iterating the proper
middle-point correspondence gives actual coherent
isomorphisms
\begin{equation}\label{eq:properclosedkernelpowers}
 \mathcal B^{\star r}
   \simeq R\Gamma(\mathbb P^1,E)^{\otimes(r-1)}
                  \otimes\mathcal B
   \simeq R\Gamma((\mathbb P^1)^{r-1},E)
                  \otimes\mathcal B
       \qquad(r\ge1).
\end{equation}
All parenthesizations integrate the same independent middle
factors, and the common product correspondence supplies
the associativity maps. The cohomological coefficients
are positive Tate multiplicities with the usual projective
Lefschetz maps. These are the classical proper product and
projective-space constructions, with the compact-support,
projection-formula and product framework in
\cite[Theorem~17.10, pp.~116--117; Proposition~18.4,
pp.~121--122; Proposition~22.5, pp.~136--137;
\S23, pp.~138--140]{Milne}.

This is a valid restricted positive model, not a new general
canonical-basis theorem. It does not make \(\mathsf G\) or
\(\mathsf P\) monoidal on the original open-kernel category.
No assertion is made that \(\mathsf G\) is naturally the
identity on a full subcategory containing every possible
derived morphism between these split objects. The actual
proper operations \eqref{eq:closedflagkernelrelation} and
\eqref{eq:properclosedkernelpowers} are the construction
being asserted.

\paragraph{Specialization and the remaining geometric scope.}
Weight grading preserves the derived Euler class and hence
any formal character and its evaluation at \(v=1\), wherever
the original comparison defines that evaluation. It still
fails to preserve the actual convolution. Perverse or ordinary
cohomology splitting removes the Chern extension and fails
the same operation comparison. Together with \fid{F62}'s
failure of tensor inversion for the torus coefficient, this
rules out these specified extraction procedures with the
unchanged standard action.

It does not rule out every filtered monoidal enhancement,
Rees construction or spectral-sequence model. Such a model
must retain the missing extension or differential data, or
change the target operation, and prove its comparison.
Equation \eqref{eq:weightgradednotmonoidal} specifies the
two objects that an unchanged-operation claim must identify.
For IV's designated nonstandard cell refinement, equality
of Euler specialization also does not identify the coordinate
submodules and their maps. If a further extraction supplied
strictly effective generator matrices in that same regular
cellular comparison, \fid{F58} would apply. No IV/VIII basis
identification is inferred from the present standard flag
example.

% Separate paragraph for Part III.
\paragraph{The weight-extraction test in Grothendieck's language.}
BBD's actual weight filtration is canonical and strictly
functorial; it produces graded perverse objects. The source
does not assert that the reconstructed additive splitting
\eqref{eq:weightgradedoperation} preserves arbitrary
convolution. Here both sides are computed using actual
six-operation correspondences and their perverse degrees.
The effective/virtual character distinction uses the
diagonalizable-character and coefficient-operation
frameworks of \cite[Expos\'e~I, 4.7.2--3]{SGA3} and
\cite[Expos\'e~V, \S\S2--3]{SGA6}, with archive anchors
\cite[cote~90, scan~17; cote~133, scan~18]{Archive}.
The zero Euler pair \(E_O\oplus E_O[-1]\) is not a zero
derived object. \(K_0\) equality therefore cannot identify
the two geometric compositions. Bounded coefficient
complexes use \cite[Expos\'e~I, Definition~4.7,
pp.~121--122, and Example~4.8, p.~122]{SGA6}; no
\(\lambda\)-ring structure is asserted for the kernel algebra.
The affine module/sheaf framework remains
\cite[Theorem~1.3.7, p.~86; Corollary~1.3.8, p.~87;
Theorem~1.4.1, pp.~90--92]{EGA}. The proper closed-kernel
model is an actual product pushforward, and its projective
coefficient follows from the usual product and
projective-space facts. No archival statement is claimed
to identify this construction with the designated
nonstandard canonical category.
